\documentclass[a4paper,12pt]{article}

\usepackage{../../../utils/template}

\newgeometry{ left=0.1cm, right=0.5cm, top=0.5cm, bottom=0.5cm }

\begin{document}
    \pagestyle{empty}

    \begin{minipage}[t]{0.46\textwidth}
        \textbf{变力沿直线做功}
        \begin{enumerate}[label=(\arabic*), leftmargin=2em, itemsep=1.2em]
            \item 设方向沿 $x$ 轴正向的力函数为 $F(x)$（$a \le x \le b$），则物体沿 $x$ 轴从点 $a$ 移动到点 $b$ 时，变力 $F(x)$ 所做的功：\\[1.5em] $W = $\vspace{2.2em}
            \item 功的元素（在 $[x,\ x+\mathrm{d}x]$ 上，$F(x)$ 近似看作常力）：\\[1.5em] $\mathrm{d}W = $\vspace{2.2em}
        \end{enumerate}

        \textbf{变速直线运动的位移}
        \begin{enumerate}[label=(\arabic*), leftmargin=2em, itemsep=1.2em]
            \item 设一物体在 $t$ 时刻的运动速度为 $v(t)$，则从 $t_1$ 时刻到 $t_2$ 时刻所经过的路程：\\[1.5em] $s = $\vspace{2.2em}
        \end{enumerate}

        \textbf{抽水做功}
        \begin{enumerate}[label=(\arabic*), leftmargin=2em, itemsep=1.2em]
            \item 将容器中的水全部抽出（$x$ 表示水面到该层水的距离，$x$ 从 $a$ 到 $b$，$\rho$ 为水的密度，$g$ 为重力加速度）：\\[1.5em] $W = $\vspace{2.2em}
            \item 功的元素（位于 $x$ 处厚度为 $\mathrm{d}x$、水平截面面积为 $A(x)$ 的一层水被抽出，路程为 $x$）：\\[1.5em] $\mathrm{d}W = $\vspace{2.2em}
        \end{enumerate}
    \end{minipage}
    \hfill
    \begin{minipage}[t]{0.46\textwidth}

        \textbf{静水压力}
        \begin{enumerate}[label=(\arabic*), leftmargin=2em, itemsep=1.2em]
            \item 垂直淹没在水中的平板 $ABCD$（水深 $x$ 从 $a$ 到 $b$）的一侧受到的水压力（$\rho$ 为水的密度，$g$ 为重力加速度）：\\[1.5em] $P = $\vspace{2.2em}
            \item 压力元素（$x$ 表示水深，$f(x) - h(x)$ 是矩形条的宽度，$\mathrm{d}x$ 是矩形条的高度）：\\[1.5em] $\mathrm{d}P = $\vspace{2.2em}
        \end{enumerate}

        \textbf{引力}
        \begin{enumerate}[label=(\arabic*), leftmargin=2em, itemsep=1.2em]
            \item 设有一长度为 $l$、线密度为常数 $\mu$ 的细棒，在细棒右端延长线距离为 $a$ 处有一质量为 $m$ 的质点 $M$，已知引力常量为 $G$，则质点 $M$ 与细棒之间的引力的大小：\\[1.5em] $F = $\vspace{2.2em}
            \item 引力元素（细棒位于 $x$ 轴上 $[-l,\ 0]$，质点位于 $x = a$ 处）：\\[1.5em] $\mathrm{d}F = $\vspace{2.2em}
        \end{enumerate}
    \end{minipage}
    \hfill

\end{document}
